\partsection{A}{The localization and the formal arguments}

\begin{claim}{A.1}{Proved}{The category and its localization}
There is a category $\Cathot_\U$ with the same objects as $\Cat_\U$ and
\[
 \Hom_{\Cathot}(X,K)=\pi _0\Fun(X,K).
\]
Here two functors represent the same morphism precisely when they are joined by a finite zigzag of natural transformations, whose orientations may vary. The class $\overline\W_\infty$ consists of the classes of functors in $\W_\infty$. The functor
$\varphi:\Cathot\to\Hot$, $[f]\mapsto[Bf]$, is a localization at this class, and
$\varphi[f]$ is invertible precisely when $f\in\W_\infty$. The identification
$\Hot\simeq\Ho(\Spc_\U)$ uses weak homotopy types of spaces, including the empty space.
\end{claim}
\begin{proof}
Whiskering a natural transformation by functors preserves its source, target, and orientation. It therefore preserves finite zigzags and defines associative composition on the displayed sets. A natural transformation $f\Rightarrow g$ is a functor $X\times[1]\to K$ restricting to $f$ and $g$. Its nerve gives a homotopy $Bf\simeq Bg$. Thus $\W_\infty$ is constant on these classes.

More intrinsically, the projection $p:X\times[1]\to X$ is in $\W_\infty$ and its two sections $i_0,i_1:X\to X\times[1]$ become the same inverse of $p$ in every category receiving a functor that inverts $\W_\infty$. Such a functor therefore identifies $f=Hi_0$ and $g=Hi_1$. The universal property of localization now factors through $\Cathot$ and gives the claimed localization. The nerve induces the equivalence with the homotopy category of simplicial sets, for example by the category-of-simplices resolution; alternatively this follows from the Quillen equivalence in A.5. Weak equivalences of simplicial sets are saturated, so $Bf$ is an isomorphism of homotopy types if and only if it is a weak equivalence. These arguments also cover $X=\varnothing$; its functor category into any $K$ is terminal.
\end{proof}

\begin{claim}{A.2}{Proved}{The specified fundamental-localizer convention}
The class $\W_\infty$ satisfies Loc(1)--(3) as stated on the first page. For any fundamental localizer in that convention, localization identifies naturally transformed functors. In particular, the categorical quotient used in A.1 is compatible with the variance of those axioms.
\end{claim}
\begin{proof}
Two-out-of-three and the stated retraction axiom hold for weak homotopy equivalences. A category with a terminal object has a natural transformation from its identity to the constant functor at that object, so its nerve is contractible. For the relative assertion, the covariant diagrams $c\mapsto A/c$ and $c\mapsto B/c$ on $C$ have Grothendieck constructions whose nerves are weakly equivalent to $NA$ and $NB$, respectively. To see the first comparison explicitly, the total category has objects $(a,c,vu(a)\to c)$. Its projection to $A$ has the section $a\mapsto(a,vu(a),\id)$, and the defining arrow supplies a natural transformation from section--projection to the identity. The construction for $B$ is identical. Objectwise weak equivalences $A/c\to B/c$ give a weak equivalence of these total categories by the homotopy-colimit theorem \citep[Theorem~26.8 and Corollary~26.9]{ChacholskiScherer2001}. This proves the relative Theorem A axiom; compare also \citet[\S1, Theorem~A]{Quillen1973}.

For an arbitrary $\W$ satisfying the stated axioms, $X/x$ and $(X\times[1])/x=(X/x)\times[1]$ have terminal objects. Their projections to the point belong to $\W$, and two-out-of-three shows that the projection of these slices belongs to $\W$. Loc(3) gives $X\times[1]\to X$ in $\W$. The section argument in A.1 then identifies naturally transformed functors in $\Cat[\W^{-1}]$. We have used the slices $A/c$, not the coslices $c\backslash A$; replacing one by the other without reversing the convention would change the hypothesis.
\end{proof}

\begin{claim}{A.3}{Proved}{Proposition 1 is a conditional formal theorem}
Let $L:C\to D$ and $R:D\to C$ be functors between locally small categories, let $L\dashv R$, and suppose $R$ is fully faithful. Write $\eta:\id_C\to RL$ and $\epsilon:LR\to\id_D$ for the unit and counit, and $W_L=\{u\mid Lu\text{ is invertible}\}$. For $K\in C$, the following conditions are equivalent: $\Hom_C(-,K)$ inverts $W_L$; it inverts every $\eta_X$; precomposition with $\eta_K$ is bijective; $\eta_K$ has a retraction; $\eta_K$ is invertible; $K$ is in the essential image of $R$; and the canonical map
\[
 \Hom_C(X,K)\longrightarrow\Hom_D(LX,LK)
\]
is bijective for every $X\in C$. If $H':D^{\op}\to\Set$ and $H=H' L^{\op}$, then $H$ is representable if and only if $H'$ is representable. In that event their representing objects satisfy $K\cong RK'$ and $K'\cong LK$, compatibly with the specified representations.
\end{claim}
\begin{proof}
Full faithfulness of $R$ makes $\epsilon$ invertible, and the triangle identity makes $L\eta_X$ invertible. The first three conditions successively imply the next, since lifting $\id_K$ under precomposition gives $p:RLK\to K$ with $p\eta_K=\id_K$. Naturality of $\eta$ gives
\[
 \eta_Kp=RL(p)\eta_{RLK}.
\]
Applying $L$ to $p\eta_K=\id$ shows that $Lp=(L\eta_K)^{-1}$. The other triangle identity and invertibility of $\epsilon$ identify $\eta_{RLK}$ with $RL\eta_K$. The right side is therefore the identity. Thus $p$ is the inverse of $\eta_K$.

If $K=RD$ up to isomorphism, adjunction and the counit identify the canonical comparison with a bijection. Conversely, that comparison, followed by adjunction, is precisely $\Hom_C(X,\eta_K)$; ordinary Yoneda makes $\eta_K$ invertible. Finally the comparison makes $\Hom_C(-,K)$ invert $W_L$. This closes the equivalences, including the apparently weaker retraction condition. If $H'$ is represented by $K'$, adjunction represents $H$ by $RK'$. If $H$ is represented by $K$, it inverts $W_L$, so the equivalence just proved identifies it with $\Hom_D(L-,LK)$. Restriction along a localization is fully faithful on presheaves; in this reflective case $L$ is such a localization. Hence $H'$ is represented by $LK$. Every uniqueness asserted here is the uniqueness supplied by Yoneda relative to the given natural isomorphism.
\end{proof}

\begin{claim}{A.4}{Proved}{Ind-categories and Proposition 2}
Use one fixed auxiliary universe for the filtered indexing categories. For $I,J$ small filtered, $K:I\to C$ and $T:J\to C$, put
\[
 h_K(X)=\colim_{i\in I}\Hom_C(X,K_i),\qquad
 \Hom_{\Ind C}(K,T)=\lim_{i\in I}\colim_{j\in J}\Hom_C(K_i,T_j).
\]
The association $K\mapsto h_K$ is fully faithful. Under the hypotheses of A.3, $\Ind L\dashv\Ind R$ and $\Ind R$ is fully faithful. For an ind-object $K$, the following are equivalent: $h_K$ inverts $W_L$; it inverts all $\eta_X$ for $X\in C$; the unit $K\to(\Ind R)(\Ind L)K$ is invertible; $K$ belongs to the essential image of $\Ind R$; and the canonical comparison
\[
 h_K(X)\longrightarrow
 \Hom_{\Ind D}(LX,(\Ind L)K)
\]
is bijective for every $X\in C$. These are precisely the valid, typed contents of V ms.~40--42.
\end{claim}
\begin{proof}
Yoneda and the universal property of the indicated colimit give
\[
 \operatorname{Nat}(h_K,h_T)
 =\lim_i\operatorname{Nat}(h_{K_i},h_T)
 =\lim_i\colim_j\Hom_C(K_i,T_j).
\]
No interchange of a limit and a colimit has occurred. Applying the adjunction inside this formula gives the ind-adjunction; applying full faithfulness inside it proves full faithfulness of $\Ind R$.

We verify the step on which the manuscript comments. Extend $h_K$ to ind-objects by $\overline h_K(X)=\Hom_{\Ind C}(X,K)=\lim_jh_K(X_j)$. If $h_K(\eta_A)$ is bijective for every $A\in C$, then $\overline h_K$ inverts the levelwise unit of every ind-object, since a limit of componentwise bijections is a bijection. Apply A.3 in $\Ind C$ to obtain the invertible unit of $K$. All remaining implications follow from A.3. Conversely, the comparison on constant source objects is exactly a natural isomorphism of the associated presheaves, and full faithfulness just proved detects the unit isomorphism. This argument uses $R$ where stated; it supplies no $R$ for $\varphi$.
\end{proof}

\begin{claim}{A.5}{Known}{The precise model-category mapping formula}
Thomason's model category has weak equivalences $\W_\infty$ and fibrations detected by $\operatorname{Ex}^2N$. Its adjunction
\[
 c\operatorname{Sd}^2:\mathbf{sSet}\rightleftarrows\Cat:
 \operatorname{Ex}^2N
\]
is a Quillen equivalence. Choose a fibrant replacement $K\to RK$ in this model category. Put $QX=c\operatorname{Sd}^2NX$. The composite last-vertex map gives a natural weak equivalence $q_X:QX\to X$, and $QX$ is cofibrant. There are natural bijections
\[
 [BX,BK]\cong\pi_0\Map_{\mathbf{sSet}}(NX,\operatorname{Ex}^2NRK)
 \cong\Fun(QX,RK)_0/\sim_T .
\]
Here $\Fun(QX,RK)_0$ is the set of functors, and $\sim_T$ is generated by endpoint restrictions of functors
$c\operatorname{Sd}^2(NX\times\Delta^1)\to RK$.
For any cofibrant replacement $Q'X\to X$, one may equivalently use $\Hom_{\Cat}(Q'X,RK)$ modulo model-category homotopy. The displayed explicit relation specifies that homotopy; it does not silently replace it by a different cylinder relation.
\end{claim}
\begin{proof}
The model structure and properness are \citet[Theorem~4.9 and Corollary~5.5]{Thomason1980}; the adjoint comparison is given in the introduction and \S5, and is recalled with the corrected cofibration argument by \citet[\S2]{Raptis2010}. Both subdivision maps and the adjunction comparison are weak equivalences, which proves the assertion about $q_X$. Since every simplicial set is cofibrant, the left Quillen functor takes $NX$ to a cofibrant category. The precise cosimplicial resolution $c\operatorname{Sd}^2(NX\times\Delta^\bullet)$ is also recorded in \citet[Remark~4.9]{Raptis2010}. The precise cosimplicial resolution $c\operatorname{Sd}^2(NX\times\Delta^\bullet)$ is also recorded in \citet[Remark~4.9]{Raptis2010}. The right-hand simplicial target is Kan by the definition of Thomason fibrancy. Its function complex therefore computes the derived mapping space out of $NX$. By adjunction, its vertices and edges are exactly the functors specified in the assertion. Connected components are the quotient by those edges. This proves the formula with its exact comparison and exact equivalence relation. The general cofibrant--fibrant formula is the defining computation of the homotopy category of a model category. In particular, calling $RK$ fibrant is not a proof of a formula with the unreplaced source $X$.
\end{proof}

\begin{claim}{A.6}{Refuted}{A fibrant target does not suffice}
The assertion that $\pi_0\Fun(X,P)\to[BX,BP]$ is always a bijection for a Thomason-fibrant $P$ is false, even for a finite $X$. The assertion that $\Fun(X,P)$ is always Thomason-fibrant is also false, even for a finite $X$ and a countable fibrant $P$.
\end{claim}
\begin{proof}
For the first assertion, let $X=B C_2$ be the one-object, two-arrow category of the cyclic group of order two. By \citet[Theorem~2.6]{Raptis2010}, choose a fibrant replacement within posets of a poset having classifying space $BC_2$. Denote it by $P$. A functor $BC_2\to P$ has to send every group element to the identity of its chosen object. Natural transformations are precisely the order relations in $P$. Thus $\Fun(BC_2,P)\cong P$. Its component set is a singleton, whereas $[BC_2,BP]$ contains a constant map and a homotopy equivalence. They induce respectively the trivial map and an isomorphism on fundamental groups, so cannot be unpointedly homotopic. The source is not cofibrant: every cofibrant category is a poset \citep[Proposition~5.7]{Thomason1980}.

For the second assertion we verify the separate construction in the companion. Let $D$ be the four-element crown, with minimal elements $a_0,a_1$, maximal elements $b_0,b_1$, and exactly the four strict relations $a_i<b_j$. Its nerve is a circle. It is not fibrant: $\operatorname{Ex}^2ND$ has finitely many edges, but a Kan complex represents each element of its based fundamental group by a single loop edge. Its fundamental group would consequently be finite, contradicting $\pi_1BD=\Z$.

Let $G=C_2$ act trivially on $D$. Starting with $Q_0=D$, at stage $r$ attach, for every Thomason generating horn $A\hookrightarrow B$ and every $F:A\to Q_r$, the free $G$-cell $G\times B$ along $(g,a)\mapsto gF(a)$. These are pushouts of Dwyer inclusions of finite posets. The old poset remains full and every new object is in a free $G$-orbit; this is the explicit pushout description in \citet[Lemmas~4.7--4.8]{MayStephanZakharevich2017}. Hence $Q=\colim_rQ_r$ is countable, fibrant, and $Q^G=D$. Fibrancy follows because each horn map has finite domain and receives a filler at a subsequent stage.

Define $E=Q\rtimes G$: an arrow $x\to y$ labelled $g$ exists when $gx\leq y$, and the composite of labels $g,h$ is $hg$. For a horn diagram $A\to E$, its labels define $A\to BG$. The horn's classifying space is contractible, so choose labels $h_a$ with $h_bg_{ab}=h_a$. Replacing the object at $a$ by $h_a x_a$ gives a diagram $A\to Q$. Extend it to $B\to Q$, extend the choices $h_a$ to the new objects, and undo the change. This produces a strict extension to $E$, proving that $E$ is fibrant.

A functor $BG\to E$ has label homomorphism either $0$ or $\id_G$. In the latter case its object satisfies $gx\leq x$, hence $gx=x$, and so lies in $D$. Naturality of a transformation labelled $h$ requires $h\phi(g)=\psi(g)h$; because $G$ is abelian it separates the two cases. Consequently
\[
 \Fun(BG,E)\cong E\amalg(D\times BG).
\]
The category $D$ is a retract of this category: project the second summand to $D$, send the first constantly to $a_0$, and use identity labels for the section. Fibrant objects are closed under retracts. Since $D$ is not fibrant, this functor category is not fibrant. The finite crown and its nerve are checked in \texttt{code/a\_finite\_checks.py}; the transfinite construction is proved above, not certified by a finite calculation.
\end{proof}

\begin{claim}{A.7}{Proved}{Descent of ind-representability needs no right adjoint}
Let $L:C\to D$ be any localization, and let $H=\colim_{i\in I}h_{K_i}$ be a small filtered colimit of representable presheaves. If $H\cong H' L^{\op}$, then the canonical comparison identifies
\[
 H'\cong\colim_{i\in I}\Hom_D(-,LK_i).
\]
For a single representing object, it follows that $h_K$ inverts the localized class if and only if its canonical comparison with $\Hom_D(L-,LK)$ is bijective. This assertion, and not a reflection, is sufficient for the descent of the manuscript's external Hom ind-objects.
\end{claim}
\begin{proof}
Restriction $L^*$ on presheaves is fully faithful, by the universal property of a localization including natural transformations. For any presheaf $F:D^{\op}\to\Set$, put $T=\colim_i h_{LK_i}$ and compute
\[
 \operatorname{Nat}(T,F)=\lim_iF(LK_i)
 =\operatorname{Nat}(H,L^*F)
 =\operatorname{Nat}(H',F).
\]
Yoneda in the presheaf category identifies $T$ with $H'$. The comparison is the one induced by $L$: this follows by tracing the identity transformation under the displayed bijections. The single-object case gives the last assertion. Thus the reverse implication, from representability on $D$ to representability on $C$, is exactly the step that is absent without $R$; the forward descent is formal.
\end{proof}
